\section{Conclusão}
\label{sec:conclusao}

O trabalho produziu uma implementação integrada de conceitos clássicos de sistemas distribuídos em uma pilha pequena e reproduzível. Os três nós executam processos independentes, trocam mensagens HTTP/JSON e mantêm estado exclusivamente local. Sobre essa base foram construídos o relógio lógico de Lamport, a exclusão mútua de Ricart--Agrawala e uma adaptação do algoritmo Bully.

Do ponto de vista de engenharia, os principais resultados são:

\begin{itemize}
    \item um envelope comum para mensagens de aplicação, mutex e eleição;
    \item atualização consistente do relógio em eventos locais, envios e recebimentos;
    \item separação explícita entre eventos Lamport e anotações de observabilidade;
    \item prioridade determinística de pedidos por timestamp e ID;
    \item observação externa da seção crítica sem centralizar a autorização;
    \item eleição, detecção de parada e recuperação do maior identificador;
    \item rejeição de coordenadores e heartbeats inferiores ao próprio nó ou ao líder conhecido;
    \item testes unitários, smoke tests e experimentos com artefatos versionados.
\end{itemize}

O relatório também esclarece duas interpretações incorretas. Primeiro, o histórico de \code{/events} não é uma sequência um-para-um de eventos teóricos: anotações podem compartilhar o timestamp do acontecimento que descrevem. Segundo, os timeouts empregados no mutex e na eleição são hipóteses operacionais da demonstração, não garantias de tolerância a falhas em um sistema assíncrono arbitrário.

Os resultados experimentais são coerentes com os comportamentos esperados: zero violações causais no cenário monitorado, quatro mensagens para um pedido individual de mutex com três nós, nenhuma sobreposição observada na seção crítica e convergência do líder após parada e recuperação.

\subsection{Limitações e possíveis extensões}

O sistema continua limitado a participantes fixos, estado volátil, rede local e falha por parada. Não há solução para partições, consenso, replicação ou persistência. Como extensões futuras, seria possível:

\begin{itemize}
    \item introduzir testes de caos para atraso, perda e reordenação de mensagens;
    \item reutilizar um \code{httpx.AsyncClient} de longa duração para pooling de conexões;
    \item ampliar os experimentos e reportar intervalos de confiança;
    \item adicionar identificadores de época mais fortes ao subsistema de eleição;
    \item comparar o Bully com um protocolo de consenso ou uma eleição baseada em log, mantendo clara a mudança de escopo;
    \item adicionar persistência apenas para fins de recuperação, sem confundi-la com o estado compartilhado proibido no núcleo da atividade.
\end{itemize}

Dentro do escopo proposto, a implementação oferece uma demonstração clara da relação entre teoria, mensagens reais, estado local e comportamento observável. O conjunto formado por código, testes, scripts, documentação e resultados permite reproduzir os cenários e discutir tanto as propriedades alcançadas quanto as garantias que permanecem fora do modelo.
